\begin{proposition}
\label{proposition-compact-is-perfect}
Let $X$ be a quasi-compact and quasi-separated scheme.
An object of $D_\QCoh(\mathcal{O}_X)$ is compact
if and only if it is perfect.
\end{proposition}

\begin{proof}
If $K$ is a perfect object of $D(\mathcal{O}_X)$ with dual
$K^\vee$ (Cohomology, Lemma \ref{cohomology-lemma-dual-perfect-complex})
we have
$$
\Hom_{D(\mathcal{O}_X)}(K, M) =
H^0(X, K^\vee \otimes_{\mathcal{O}_X}^\mathbf{L} M)
$$
functorially in $M$. Since $K^\vee \otimes_{\mathcal{O}_X}^\mathbf{L} -$
commutes with direct sums and since $H^0(X, -)$ commutes with direct
sums on $D_\QCoh(\mathcal{O}_X)$ by
Lemma \ref{lemma-quasi-coherence-pushforward-direct-sums}
we conclude that $K$ is compact in $D_\QCoh(\mathcal{O}_X)$.

\medskip\noindent
Conversely, let $K$ be a compact object of $D_\QCoh(\mathcal{O}_X)$.
To show that $K$ is perfect, it suffices to show that
$K|_U$ is perfect for every affine open $U \subset X$, see
Cohomology, Lemma \ref{cohomology-lemma-perfect-independent-representative}.
Observe that $j : U \to X$ is a quasi-compact and separated morphism.
Hence
$Rj_* : D_\QCoh(\mathcal{O}_U) \to D_\QCoh(\mathcal{O}_X)$
commutes with direct sums, see
Lemma \ref{lemma-quasi-coherence-pushforward-direct-sums}.
Thus the adjointness of restriction to $U$ and $Rj_*$ implies that
$K|_U$ is a compact object of $D_\QCoh(\mathcal{O}_U)$.
Hence we reduce to the case that $X$ is affine.

\medskip\noindent
Assume $X = \Spec(A)$ is affine. By Lemma \ref{lemma-affine-compare-bounded}
the problem is translated into the same problem for $D(A)$.
For $D(A)$ the result is
More on Algebra, Proposition \ref{more-algebra-proposition-perfect-is-compact}.
\end{proof}
